\documentclass[a4paper,12pt]{article}

\usepackage[utf8]{inputenc}
\usepackage[T1]{fontenc}
\usepackage[ngerman]{babel}
\usepackage{lmodern}
\usepackage{amsmath}
\usepackage{graphicx}
\usepackage{hyperref}
\usepackage{array}

\title{LaTeX-Testdokument für TUI-Editoren}
\author{Dr. Michael H. Raus}
\date{\today}

\begin{document}

\maketitle
\tableofcontents

\section{Überschriften und Abschnitte}

Dies ist ein Testabschnitt mit normalem Fließtext. Er enthält mehrere Sätze,
damit Zeilenumbrüche, Worttrennung und Absatzverhalten geprüft werden können.
Der Text dient ausschließlich dazu, typische LaTeX-Strukturen sichtbar zu
machen.

\subsection{Unterabschnitt}

Auch Unterabschnitte sollen erkannt werden. Dieser Absatz enthält deutsche
Umlaute wie ä, ö, ü und ß. Außerdem enthält er Satzzeichen, Klammern
(rund), [eckig] und \{geschweift\}.

\subsubsection{Unter-Unterabschnitt}

Dies ist eine tiefere Gliederungsebene. Sie eignet sich zum Testen von
Syntax-Hervorhebung, Folding und Navigationsfunktionen im Editor.

\section{Absätze und Zeilenumbrüche}

Dies ist der erste Absatz. Ein Absatz endet in LaTeX durch eine Leerzeile.

Dies ist der zweite Absatz.
Hier wurde ein manueller Zeilenumbruch mit zwei Leerzeichen am Zeilenende
simuliert, obwohl LaTeX normalerweise eigene Zeilenumbrüche setzt.

Eine neue Zeile kann auch explizit erzeugt werden.\\
Dies ist die Zeile nach einem expliziten Umbruch.

\section{Textauszeichnungen}

LaTeX unterstützt verschiedene Textauszeichnungen. Hier ist ein Wort
\textbf{fett}, ein anderes \textit{kursiv}, ein weiteres \underline{unterstrichen}
und dieses hier \emph{betont}.

Man kann Auszeichnungen auch kombinieren:
\textbf{\textit{fett und kursiv}}.

Auch Schreibmaschinenschrift ist möglich:
\texttt{Dies ist Monospace-Text.}

\section{Listen}

\subsection{Ungeordnete Liste}

\begin{itemize}
    \item Erster Punkt einer Liste.
    \item Zweiter Punkt mit etwas längerem Text, damit Umbrüche getestet werden.
    \item Dritter Punkt mit \textbf{fetter Hervorhebung}.
\end{itemize}

\subsection{Geordnete Liste}

\begin{enumerate}
    \item Erster nummerierter Punkt.
    \item Zweiter nummerierter Punkt.
    \item Dritter nummerierter Punkt.
\end{enumerate}

\subsection{Beschreibungsliste}

\begin{description}
    \item[Name] Ein beschreibender Eintrag.
    \item[Typ] Eine weitere Beschreibung.
    \item[Zweck] Test von Einrückung und Label-Darstellung.
\end{description}

\section{Zitate und Umgebungen}

Ein kurzes Zitat kann mit der Umgebung \texttt{quote} gesetzt werden.

\begin{quote}
Dies ist ein eingerücktes Zitat. Es besteht aus mehreren Sätzen und prüft,
ob der Editor Umgebungen korrekt erkennt.
\end{quote}

Eine längere Passage kann auch in einer eigenen Umgebung stehen:

\begin{quotation}
Dies ist eine längere Zitierumgebung. Sie kann mehrere Absätze enthalten.

Dies ist ein zweiter Absatz innerhalb derselben Umgebung.
\end{quotation}

\section{Fußnoten und Verweise}

Dies ist ein Satz mit einer Fußnote.\footnote{Dies ist der Text der Fußnote.}

Auf Abschnitt \ref{sec:mathematik} wird später verwiesen. Hyperlinks werden
durch das Paket \texttt{hyperref} erzeugt.

\section{Mathematik}
\label{sec:mathematik}

Inline-Mathematik steht zwischen Dollarzeichen, zum Beispiel $a^2 + b^2 = c^2$.

Abgesetzte Mathematik kann so aussehen:

\[
    E = mc^2
\]

Eine nummerierte Gleichung:

\begin{equation}
    \int_0^1 x^2 \, dx = \frac{1}{3}
\end{equation}

Eine mehrzeilige Gleichungsumgebung:

\begin{align}
    f(x) &= x^2 + 2x + 1 \\
         &= (x + 1)^2
\end{align}

\section{Tabellen}

Eine einfache Tabelle:

\begin{center}
\begin{tabular}{|l|c|r|}
\hline
Name & Anzahl & Preis \\
\hline
Apfel & 3 & 1,20 EUR \\
Birne & 5 & 2,50 EUR \\
Banane & 12 & 4,80 EUR \\
\hline
\end{tabular}
\end{center}

Eine Tabelle mit fester Spaltenbreite:

\begin{center}
\begin{tabular}{|p{4cm}|p{8cm}|}
\hline
Begriff & Beschreibung \\
\hline
Editor & Ein Programm zum Bearbeiten von Textdateien. \\
\hline
Syntax & Die formale Struktur einer Sprache oder eines Dateiformats. \\
\hline
\end{tabular}
\end{center}

\section{Abbildungen}

Hier folgt eine Beispielabbildung. Die Datei muss nicht existieren, wenn nur
die Syntax des Dokuments getestet werden soll. Für echtes Kompilieren muss
\texttt{example-image} verfügbar sein, etwa über das Paket \texttt{mwe}.

\begin{figure}[h]
    \centering
    \fbox{\rule{0pt}{4cm}\rule{6cm}{0pt}}
    \caption{Platzhalter für eine Abbildung}
    \label{fig:platzhalter}
\end{figure}

Abbildung \ref{fig:platzhalter} zeigt einen einfachen Platzhalterrahmen.

\section{Quellcode und Verbatim}

Kurzer Inline-Code kann mit \verb|printf("Hallo Welt\n");| geschrieben werden.

Eine Verbatim-Umgebung:

\begin{verbatim}
#include <stdio.h>

int main(void) {
    printf("Hallo Welt\n");
    return 0;
}
\end{verbatim}

In der Verbatim-Umgebung bleiben Sonderzeichen wie \#, \%, \&, \_, \{ und \}
unverändert sichtbar.

\section{Sonderzeichen und Kommentare}

Einige Zeichen haben in LaTeX besondere Bedeutung und müssen maskiert werden:

\begin{itemize}
    \item Prozentzeichen: \%
    \item Kaufmännisches Und: \&
    \item Unterstrich: \_
    \item Dollarzeichen: \$
    \item Raute: \#
    \item Geschweifte Klammern: \{ und \}
    \item Backslash: \textbackslash
\end{itemize}

% Dies ist ein Kommentar. Er erscheint nicht im gesetzten Dokument.

Dieser Satz steht nach einem Kommentar und sollte normal verarbeitet werden.

\end{document}
